\begin{figure}[ht]
\centering
\begin{tikzpicture}[x=1.1cm, y=11cm, font=\small]

% Partial sums of alternating harmonic series converging to log 2 = 0.6931
% S1=1.0000, S2=0.5000, S3=0.8333, S4=0.5833, S5=0.7833, S6=0.6167,
% S7=0.7595, S8=0.6345, S9=0.7456, S10=0.6456
% y range ~ [0.45, 1.05]; subtract 0.45 baseline, scale handled by y unit.
\def\base{0.45}

% Axes
\draw[->, thick] (0.5, 0) -- (10.7, 0) node[right] {$N$};
\draw[->, thick] (0.5, 0) -- (0.5, 0.62) node[above] {$S_{N}$};

% y ticks at 0.5, 0.6931, 0.8, 1.0 (shifted by base)
\foreach \v/\lbl in {0.5/0.5, 0.6931/$\log 2$, 0.8/0.8, 1.0/1.0} {
  \pgfmathsetmacro\yy{\v-\base}
  \draw (0.45, \yy) -- (0.55, \yy) node[left, xshift=-0.1cm] {\lbl};
}
% x ticks
\foreach \n in {1,2,...,10} {
  \draw (\n, -0.005) -- (\n, 0.005);
}
\node[below] at (2, -0.01) {2};
\node[below] at (4, -0.01) {4};
\node[below] at (6, -0.01) {6};
\node[below] at (8, -0.01) {8};
\node[below] at (10, -0.01) {10};

% limit line
\draw[dashed, gray] (0.5, {0.6931-\base}) -- (10.6, {0.6931-\base});

% odd partial sums (approach from above) - red, decreasing
\draw[thick, engred, mark=*, mark size=1pt]
  plot coordinates {
  (1, {1.0000-\base}) (3, {0.8333-\base}) (5, {0.7833-\base})
  (7, {0.7595-\base}) (9, {0.7456-\base}) };
\node[engred, anchor=west] at (9.1, {0.7456-\base}) {odd $S_N$};

% even partial sums (approach from below) - blue, increasing
\draw[thick, engblue, mark=square*, mark size=1pt]
  plot coordinates {
  (2, {0.5000-\base}) (4, {0.5833-\base}) (6, {0.6167-\base})
  (8, {0.6345-\base}) (10, {0.6456-\base}) };
\node[engblue, anchor=west] at (10.1, {0.6456-\base}) {even $S_N$};

\end{tikzpicture}
\caption[Nested-interval envelope of alternating partial sums]{Partial
sums of the alternating harmonic series $1 - 1/2 + 1/3 - \cdots$
converging to $\log 2$. The odd partial sums descend from above and the
even partial sums climb from below, trapping the limit in a nested
sequence of intervals $[S_{2k}, S_{2k+1}]$ whose width is the next term
$b_{2k+1}$. The width is the alternating-series error bound made
visible: any partial sum lies within one term of the limit.}
\label{fig:vol02:ch04:alternating-envelope}
\end{figure}
